\chapter{Introducción}\label{cap:intro}

\section{El problema del viajante}

El \emph{problema del viajante} (en inglés, \emph{travelling salesman problem}, TSP) es uno de los problemas clásicos de la optimización
combinatoria. Un viajante debe visitar un conjunto de ciudades, pasando por cada una exactamente una vez, y volver a la ciudad de partida
recorriendo la menor distancia posible. Aunque el enunciado es sencillo, encontrar la mejor ruta es difícil: con $n$ ciudades hay $(n-1)!$ recorridos
posibles, un número que crece más deprisa que cualquier exponencial.

Formalmente, sea $V=\{0,1,\dots,n-1\}$ el conjunto de ciudades y $c(i,j)\ge 0$ el coste de ir de la ciudad $i$ a la ciudad $j$. Una ruta es una permutación
$\pi$ de $V$ y su coste, incluida la vuelta al inicio, es
\begin{equation}\label{eq:coste}
  C(\pi)\;=\;\sum_{k=0}^{n-2} c\bigl(\pi_k,\pi_{k+1}\bigr)\;+\;c\bigl(\pi_{n-1},\pi_0\bigr).
\end{equation}
El problema consiste en encontrar la permutación $\pi^{*}$ que minimiza $C(\pi)$. Si $c(i,j)=c(j,i)$ para todo par de ciudades el problema es
\emph{simétrico}; en caso contrario es \emph{asimétrico}. Si además se cumple la \emph{desigualdad triangular},
\begin{equation}\label{eq:triangular}
  c(i,k)\;\le\;c(i,j)+c(j,k)\qquad\text{para todo }i,j,k\in V,
\end{equation}
decimos que el problema es \emph{métrico}; esta propiedad es necesaria para algunas garantías de aproximación, como la del algoritmo de Christofides
(sección~\ref{sec:christofides}).

El TSP es un problema NP-difícil~\cite{karp1972,garey1979}: no se conoce ningún algoritmo que lo resuelva en tiempo polinómico en el caso
general, y se cree que no existe. Eso no significa que sea imposible resolverlo en la práctica. Los algoritmos exactos modernos, basados en programación
entera, han resuelto instancias de decenas de miles de ciudades~\cite{applegate2006}, y las heurísticas obtienen soluciones muy buenas en
tiempos pequeños. La cuestión interesante es \emph{cuánto cuesta} cada enfoque y \emph{qué se sacrifica} a cambio de rapidez, que es lo que estudia
este trabajo.

\section{Motivación: carreteras reales}

La mayoría de los estudios sobre el TSP usan instancias sintéticas (puntos en el plano con distancia euclídea) o colecciones clásicas. Aquí lo
planteamos sobre un caso real: \emph{las ciudades de España y las carreteras que las unen}. Esto añade tres ingredientes que no aparecen en el plano:
\begin{itemize}
  \item \textbf{Las distancias son por carretera}, es decir, caminos mínimos sobre un grafo vial con cientos de miles de nodos, no distancias en línea recta.
  \item \textbf{Hay vías de sentido único} (por ejemplo, las dos calzadas de una autovía son aristas distintas), de modo que ir de $A$ a $B$ puede costar distinto que de $B$ a $A$.
        El problema es entonces asimétrico, y la elección entre un modelo dirigido y otro no dirigido pasa a ser parte del estudio.
  \item \textbf{Se puede medir en kilómetros o en tiempo}, y la ruta más corta no tiene por qué ser la más rápida.
\end{itemize}

Estos ingredientes permiten, además, plantear preguntas que un benchmark sintético no permite: ¿cómo cambia el recorrido óptimo si solo se pueden usar autovías?
¿Cuánto cuesta ignorar los sentidos únicos? ¿Qué algoritmos se pueden aplicar a qué versión del problema, y por qué?

\section{Objetivos}

El objetivo general es \textbf{medir cómo mejoran en complejidad y en tiempo los distintos algoritmos del viajante al buscar una solución óptima o completa, y qué
se pierde en calidad cuando se renuncia al óptimo}. Lo desglosamos en cinco preguntas, que guían todo el estudio:
\begin{enumerate}[label=\textbf{P\arabic*.}]
  \item ¿Qué algoritmo puede usarse con qué versión del grafo, y por qué?
  \item ¿Cómo crecen el tiempo y la memoria con el número de ciudades? ¿Coincide con la complejidad teórica?
  \item ¿Cuánto se aleja de la ruta óptima cada algoritmo que no garantiza el óptimo?
  \item ¿Cómo cambian los resultados con el tipo de carreteras, el peso (kilómetros o tiempo) y el sentido de las vías?
  \item ¿Qué algoritmo conviene en cada caso, atendiendo a calidad y tiempo?
\end{enumerate}
Para responderlas hay que resolver antes problemas previos de ingeniería de datos: obtener un grafo vial fiable, elegir y justificar un conjunto de ciudades, y
definir un conjunto de \emph{versiones} del problema bien especificadas y reproducibles.

\section{Qué se ha hecho}

El trabajo se organiza en cinco partes encadenadas (figura~\ref{fig:pipeline}), a las que se suman una web interactiva y esta memoria:
\begin{enumerate}[label=\textbf{Parte \arabic*.},leftmargin=*,labelsep=0.5em]
  \item \textbf{Construcción del grafo.} Descarga de un extracto de OpenStreetMap, filtrado, simplificación y asignación de longitudes y tiempos (capítulo~\ref{cap:grafo}).
  \item \textbf{Ciudades y versiones del grafo.} Selección de 75 ciudades y cálculo de tablas de distancias que originan 480 versiones del problema (capítulo~\ref{cap:versiones}).
  \item \textbf{Algoritmos.} Implementación de ocho algoritmos, tres exactos, tres heurísticos y dos metaheurísticos (capítulo~\ref{cap:algoritmos}).
  \item \textbf{Estudio empírico.} Ejecución sistemática de los algoritmos y análisis de tiempo, memoria y calidad (capítulos~\ref{cap:metodologia} y~\ref{cap:resultados}).
  \item \textbf{Conclusiones.} Respuesta a las preguntas planteadas y recomendaciones de uso (capítulo~\ref{cap:conclusiones}).
\end{enumerate}

\begin{figure}[htbp]
  \centering
  \begin{tikzpicture}[
      caja/.style={draw=#1!70!black,fill=#1!10,rounded corners=3pt,align=center,minimum height=1.25cm,text width=2.9cm,font=\small,line width=0.6pt},
      flecha/.style={-{Stealth[length=2.2mm]},line width=0.7pt,color=gris}]
    \node[caja=azul] (a) {\textbf{OpenStreetMap}\\extracto de España\\(\texttt{.osm.pbf}, 1,4~GB)};
    \node[caja=naranja,right=0.75cm of a] (b) {\textbf{Grafo vial}\\179.594 nodos\\309.250 aristas};
    \node[caja=aqua,right=0.75cm of b] (c) {\textbf{Versiones}\\16 tablas de distancias\\240 grafos de ciudades};
    \node[caja=violeta,below=0.95cm of c] (d) {\textbf{8 algoritmos}\\exactos, heurísticos,\\metaheurísticas};
    \node[caja=rojo,left=0.75cm of d] (e) {\textbf{Estudio empírico}\\4.706 ejecuciones\\tiempo, memoria, calidad};
    \node[caja=verde,left=0.75cm of e] (f) {\textbf{Conclusiones}\\y recomendaciones\\de uso};
    \node[caja=amarillo,below=0.9cm of e] (g) {\textbf{Web interactiva}\\animación paso a paso\\y carrera};
    \draw[flecha] (a) -- (b); \draw[flecha] (b) -- (c); \draw[flecha] (c) -- (d); \draw[flecha] (d) -- (e); \draw[flecha] (e) -- (f);
    \draw[flecha,dashed] (d) |- (g.east);
  \end{tikzpicture}
  \caption{Esquema del trabajo: de los datos de OpenStreetMap a las conclusiones. La web reutiliza los datos y los algoritmos.}
  \label{fig:pipeline}
\end{figure}

\section{Herramientas y reproducibilidad}

Todo el trabajo está escrito en Python~3 y se entrega como un \emph{notebook} (\texttt{trabajo\_tsp.ipynb}) más tres scripts de apoyo (los \emph{anexos}):
el Anexo~A descarga el extracto, el Anexo~B construye el grafo y el Anexo~C ejecuta el estudio empírico en paralelo. Usamos NetworkX~\cite{hagberg2008}
para los grafos, NumPy~\cite{harris2020} y SciPy~\cite{virtanen2020} para el cálculo numérico, PuLP con el \emph{solver} CBC para la programación entera
de DFJ y OSMnx~\cite{boeing2017} únicamente en la construcción del grafo. El código, los datos intermedios y los resultados se versionan en GitHub, y cada
resultado de esta memoria se puede regenerar con los comandos del apéndice~\ref{ap:reproducir}.

\section{Estructura de la memoria}

El capítulo~\ref{cap:grafo} describe los datos y la construcción del grafo vial. El capítulo~\ref{cap:versiones} presenta las ciudades y las versiones del problema.
El capítulo~\ref{cap:algoritmos} explica los ocho algoritmos y justifica su elección. El capítulo~\ref{cap:metodologia} detalla la metodología del estudio y el
capítulo~\ref{cap:resultados} presenta sus resultados. El capítulo~\ref{cap:conclusiones} recoge las conclusiones. El capítulo~\ref{cap:web} describe la web interactiva
y el capítulo~\ref{cap:organizacion} la organización del trabajo en equipo, el uso de herramientas de inteligencia artificial y los problemas encontrados. Los apéndices
reúnen las instrucciones para reproducir el trabajo y material complementario.
